\documentclass[11pt]{article}
\usepackage[margin=0.85in]{geometry}
\usepackage{amsmath,amssymb,booktabs}
\usepackage[hidelinks]{hyperref}
\setlength{\parindent}{0pt}
\setlength{\parskip}{0.35em}
\setlength{\emergencystretch}{2em}
\newcommand{\ve}{\mathrm{veh\mbox{-}eq}}
\newcommand{\dd}{\mathrm{d}}
\title{VSL-R1: What Makes Extra Preparation Time Useful?\\[0.25em]
\large Conditional transfer-clock and receiving audit}
\author{Method-first task-definition review}
\date{12 October 2026}
\begin{document}
\maketitle
\vspace{-1em}

\textbf{Status and purpose.} These are analytical audit deductions, with no
simulation or new primary traffic law adopted. S13 remains the accepted
conditional reference. The audit stops at gate VSL-R1-G1 because a fixed
physical-time transfer rate and a flow/distance-coupled transfer are different
models under common-speed control.

\section*{1. The marginal benefit the inherited task needs}
Keep S13's fixed affected mass $M>0$ in $\ve$, zone length $L>0$ in metres,
common realized speed $v>0$ in $\mathrm{m/s}$, receiving-work factor
$\alpha>1$ and work rate $\mu>0$ in normalized $\ve/\mathrm{s}$.
The ideal traversal time is $\tau=L/v$. Let $E(v)$ be the mass remaining in
the ending lane at the choke. Conditional on the \emph{same proportional
batch-service law}, the mean complete time is
\begin{equation}
 J(v)=\frac{L}{v}+\frac{M}{2\mu}+aE(v),
 \qquad a=\frac{\alpha-1}{2\mu}.
 \label{eq:objective}
\end{equation}
$J$ is seconds; $a$ has units $\mathrm{s}/\ve$. Total time is $F=MJ$,
in $\ve\,\mathrm{s}$. Consequently, in a smooth branch,
\begin{equation}
 J'(v)=-\frac{L}{v^2}+aE'(v).
 \label{eq:marginal}
\end{equation}
A small feasible slowdown $\delta v<0$ improves $J$ to first order only if
\begin{equation}
 J'(v)>0\quad\Longleftrightarrow\quad
 aE'(v)>\frac{L}{v^2}.
 \label{eq:criterion}
\end{equation}
The first term is the added travel time everyone pays. The second is the
reduction in mean service burden under this particular allocation. Finite
receiving and positive lateral transfers alone do not establish the inequality.
At the fastest bound the slowdown is feasible; an increase beyond the bound
is not. Equality gives no first-order conclusion.

This expression must not be reused unchanged if preparation interaction delay,
arrival chronology, receiving or allocation is altered. Those changes require
their own complete departure-time accounting. The fixed coefficient $a$ is
not a source-independent service identity.

\newpage
\section*{2. Holding a time rate fixed is a scientific assumption}
For a diagnostic, suppose a nonbinding, prescribed-motion preparation uses a
fractional completion rate $\lambda(v)>0$ in $\mathrm{s}^{-1}$, constant along
one traversal at a fixed action. With initial ending mass $E_0=pM$,
\begin{equation}
 E(v)=E_0\exp\!\left[-\lambda(v)\frac{L}{v}\right].
 \label{eq:rate}
\end{equation}
This diagnostic family is \textbf{not adopted as a new traffic model}.
It exposes what a source-to-reduction argument would have to justify. Direct
differentiation gives
\begin{align}
 E'(v)&=\frac{LE(v)}{v^2}
          \bigl[\lambda(v)-v\lambda'(v)\bigr],\label{eq:rategrad}\\
 J'(v)&=\frac{L}{v^2}
       \left[-1+aE(v)\bigl(\lambda(v)-v\lambda'(v)\bigr)\right].
 \label{eq:joinedgrad}
\end{align}
The speed dependence of the rate matters as much as its positive value.

\textbf{Inherited S13 clock.} Hold the physical-time rate fixed:
$\lambda(v)=\lambda_0>0$. Then
\begin{equation}
 E(v)=E_0e^{-\lambda_0L/v},\qquad
 J'(v)=\frac{L}{v^2}\bigl[-1+a\lambda_0E(v)\bigr].
 \label{eq:s13}
\end{equation}
This restates the accepted sensitivity and its conditional improving regime.

\textbf{Speed-scaled countercase.} Hold a per-distance completion coefficient
$\beta>0$ in $\mathrm{m}^{-1}$ fixed, so $\lambda(v)=\beta v$. Then
\begin{equation}
 E(v)=E_0e^{-\beta L},\qquad E'(v)=0,\qquad
 J'(v)=-L/v^2<0.
 \label{eq:spatial}
\end{equation}
Longer residence is exactly offset by a smaller per-time rate. Under the
inherited service law, fastest speed is best for every preparation even though
conservative merging and extra receiving work remain present.

\textbf{Source connection and limit.} FB0030's equation 14 makes desired
lateral flow a fraction of sending demand; its free-flow sending branch
(equation 13) contains speed times density. That motivates checking rate scaling.
Equation~\eqref{eq:spatial} is an original conditional diagnostic, not a
continuum derivation or VSL prediction of the full native two-step model.
Density, incentives, receiving and cooperation may break this scaling.
Carey's location-based desire likewise is not a prescribed realized-completion
schedule. Neither source identifies $\lambda'(v)=0$ for this task.

\newpage
\section*{3. Conditional receiving reduction behind batch co-clearance}
The audited Carey R2 allocation uses gap-weighted receiving demands. Here we
project that interpretation onto two simultaneously available queues, rather
than claiming to reproduce the source's spatial update system.

At the start of service let $e(0)=E$, $c(0)=M-E$. For remaining masses $c,e$,
assume constant work rate $\mu$, unit work for continuing mass and $\alpha$
work for ending mass. All demand is exposed, with no other binding constraint.
Proportional allocation to work demands $c$ and $\alpha e$ gives
\begin{equation}
 \dot c=-\frac{\mu c}{c+\alpha e},\qquad
 \dot e=-\frac{\mu e}{c+\alpha e},
 \qquad c+e>0.
 \label{eq:allocation}
\end{equation}
For positive classes their ratio is invariant; zero classes remain zero.
Writing $B=c+e$, $f=E/M$ and $\overline w=1+(\alpha-1)f$ yields
\begin{equation}
 c+\alpha e=\overline w B,\qquad
 \dot B=-\mu/\overline w,\qquad
 T=\frac{\overline w M}{\mu}=\frac{M+(\alpha-1)E}{\mu}.
 \label{eq:linear}
\end{equation}
Thus, for $0\le u\le T$,
\begin{equation}
 B(u)=M(1-u/T),\qquad
 \int_0^T B(u)\,\dd u=\frac{MT}{2}.
 \label{eq:area}
\end{equation}
Service stops at $B=0$; the quotient is not evaluated there. Work is consumed
at rate $\mu$, and discharged traffic is the removed mass, not the work.
Including the common preparation traversal gives $F=M\tau+MT/2$, as in S13.

This is a coherent conditional resource model. It requires a simultaneous
batch, constant accessible receiving, the stated allocation and no independent
donor/sending limitation. A spatial CTM generally does not present its entire
population as those two queue demands. Sharing a target cell is also different
from assuming both continuing lanes offer one accessible resource.

\textbf{Preparation receiving.} Nonbinding preparation is a separate regime
assumption. A lane change can consume receiving space without adding travel
delay when that resource has slack, but that slack must be established over the
action range. Slower motion can change density and slack. The choke projection
does not provide such a regime bound or a physical-time completion law.

\textbf{Information and gate.} $E_0$ alone still reproduces the inherited
grouped model under its fixed law. No future reference $E(v)$ is supplied.
Native receiving-limited extensions may need adjacent-lane state; none is
selected here. VSL-R1-G1 concerns retaining an explicit conditional benchmark
versus amending its primary preparation/receiving contract. Independent useful
traffic gradients, calibration and implementation checks remain unestablished.

\end{document}
